\section{M0 --- Infraestrutura e Reprodutibilidade}
\label{sec:m0}

A primeira etapa estabeleceu a infraestrutura que garante reprodutibilidade e
separação entre lógica e orquestração: a lógica reside em \texttt{src/} e os notebooks
apenas orquestram e visualizam.

\subsection{Ambiente e dependências}
Fixou-se o interpretador em \textbf{Python 3.12.13} (via \texttt{pyenv}), versão idêntica
à do Google Colab à época, alvo principal de execução. As dependências foram
\emph{pinadas} (TensorFlow/Keras~\cite{tensorflow2015}, \texttt{yfinance}, \texttt{pandas},
\texttt{numpy}, \texttt{scikit-learn}, \texttt{matplotlib}, \texttt{seaborn}), habilitando
instalação editável (\texttt{pip install -e .}). Notebooks têm suas saídas removidas
automaticamente no versionamento (\texttt{nbstripout}), reduzindo ruído em \emph{diffs}.

\subsection{Configuração centralizada e seeds}
Todos os hiperparâmetros residem em \texttt{config.yaml}, fonte única de verdade: alterar
um experimento significa alterar esse arquivo, não os notebooks. O módulo
\texttt{src/config.py} carrega essa configuração e expõe \texttt{set\_seeds()}, que fixa
as sementes de \texttt{random}, \texttt{numpy} e \texttt{tensorflow} para reprodutibilidade.

\subsection{Decisões de hiperparâmetros e proveniência}
Um princípio metodológico central do projeto é a \textbf{rastreabilidade da proveniência}
de cada hiperparâmetro. Cada escolha é classificada como: \emph{fundamentada}
(ancorada no enunciado do projeto ou na literatura), \emph{default de literatura}
(convenção da implementação de referência), \emph{decisão de projeto} (escolha de
\emph{design} sem fonte direta, com \emph{rationale} explícito) ou \emph{provisória}
(valor a ser calibrado por experimento). A Tabela~\ref{tab:hparams} consolida as
principais decisões; os valores provisórios não devem sustentar resultados finais sem
validação.

\begin{table}[h]
\centering
\caption{Principais hiperparâmetros e sua proveniência (consolidado dos ADRs).}
\label{tab:hparams}
\small
\begin{tabular}{lll}
\toprule
\textbf{Parâmetro} & \textbf{Valor} & \textbf{Proveniência} \\
\midrule
\texttt{window\_size}        & 30      & Fundamentada~\cite{li2020lstm,valkov2019lstm} \\
\texttt{lstm\_units}         & 64      & Default de literatura~\cite{valkov2019lstm} \\
\texttt{dropout}             & 0{,}2   & Default de literatura~\cite{valkov2019lstm} \\
\texttt{latent\_dim}         & 16      & Decisão de projeto (calibrar) \\
\texttt{loss}                & MSE     & Fundamentada~\cite{li2020lstm} \\
\texttt{shuffle}             & False   & Fundamentada (causalidade)~\cite{valkov2019lstm} \\
\texttt{validation\_split}   & 0{,}1   & Default de literatura~\cite{valkov2019lstm} \\
\texttt{threshold\_percentile} & 95    & Fundamentada (enunciado) \\
\texttt{dynamic\_window}     & 252     & Provisória ($\approx$1 ano útil) \\
\texttt{n\_injections}       & 50      & Provisória \\
\bottomrule
\end{tabular}
\end{table}

\subsection{Riscos metodológicos registrados}
Três riscos foram documentados já nesta fase, para serem honrados na implementação:
(i) o \texttt{validation\_split} do Keras separa o \emph{bloco final} dos dados antes de
embaralhar --- logo \texttt{shuffle=False} é obrigatório, sob pena de vazamento temporal;
(ii) o limiar dinâmico deve ser \emph{causal}, usando apenas observações passadas;
(iii) o \texttt{MinMaxScaler} ajustado no treino faz com que extremos do teste possam
exceder o intervalo $[0,1]$ --- comportamento esperado e desejável, pois é onde residem as
anomalias, mas que torna o limiar sensível a \emph{outliers} isolados (motivando o uso de
percentil em vez do máximo).
